\subsection{联合振幅怎样决定纠缠与条件光态}
\label{chap:feqo-conditioning}
求得联合振幅以后，电子谱仍只保存其中的一部分信息。
若不同电子通道留下不同光学记录，对光求迹会减小电子相干；
若再测量这些记录，可以选择条件电子态。
这两个操作都由同一系数矩阵计算。

\begin{definition}[联合纯态的 Schmidt 权重与纠缠熵]
对可分 Hilbert 空间中的归一化联合纯态
$|\Psi\rangle=\sum_{\ell,n}C_{\ell n}|\ell,n\rangle$，
令 $\rho_e=\operatorname{Tr}_{\rm ph}|\Psi\rangle\langle\Psi|$。
其非零本征值 $\lambda_k$ 称为 Schmidt 权重，$\sqrt{\lambda_k}$ 是 Schmidt 系数。
纠缠熵定义为
\begin{equation}
S_e=-\operatorname{Tr}(\rho_e\ln\rho_e)
=-\sum_k\lambda_k\ln\lambda_k,\qquad0\ln0=0.
\label{eq:feqo-entanglement-entropy}
\end{equation}
允许 $S_e=+\infty$；这仍是熵的定义，不需以任意截断值代替它。
\end{definition}

\begin{theorem}[联合振幅的奇异值与可分性]
系数矩阵 $C$ 是 Hilbert--Schmidt 算符，且存在展开
\begin{equation}
|\Psi\rangle=\sum_k\sqrt{\lambda_k}|u_k\rangle_e|v_k\rangle_{\rm ph},
\qquad\sum_k\lambda_k=1.
\label{eq:feqo-schmidt-decomposition}
\end{equation}
联合纯态可分当且仅当 $C$ 秩为一，亦即所有非零列共线。
若有限系统的较小维数为 $d$，则 $0\le S_e\le\ln d$，
上界当且仅当 $d$ 个权重全为 $1/d$。
\end{theorem}
\begin{proof}
归一化给 $\sum_{\ell,n}|C_{\ell n}|^2=1$，因此 $C$ 为紧的 Hilbert--Schmidt 算符。
其奇异值分解 $C=\sum_k s_k u_k w_k^\dagger$ 满足 $\sum_ks_k^2=1$。
把矩阵元代回联合态，光侧向量取 $|v_k\rangle=\sum_n(w_k)_n^*|n\rangle$，
即可得到 $\lambda_k=s_k^2$ 的展开。
对任一子系统求迹，另一侧正交性给相同的非零本征值。
只有一个奇异值时是积态；反向把积态系数写为 $C_{\ell n}=u_\ell v_n$ 就证明秩一。
\begin{equation}
\operatorname{Tr}\rho_e^2
=\sum_{n,m}|\langle c_n|c_m\rangle|^2\le
\sum_{n,m}\|c_n\|^2\|c_m\|^2=1,
\label{eq:feqo-entanglement-column-test}
\end{equation}
其中 $c_n$ 是第 $n$ 列；等号要求所有非零列满足 Cauchy--Schwarz 等号条件，即共线。
有限维熵上界由
$\sum_k\lambda_k\ln(d\lambda_k)\ge0$ 得到，等号要求均匀权重。
对于混合联合态，局部混合可以来自经典概率；局部熵本身不能直接称为纠缠熵。
\end{proof}

\begin{exbox}[真空输入：Poisson 交换与纠缠熵的完整计算]
取电子单格、光真空的因子化输入，求联合末态、电子纯度与纠缠熵，并判断相同的 Poisson 数分布是否足以辨认光相干态。

\textbf{解。} 令 $\zeta=-g_{\rm Q}^*$、$\mu=|g_{\rm Q}|^2$。
由位移算符的正规排序作用于真空，单模输出为
\begin{equation}
|\Psi\rangle=\ee^{-\mu/2}\sum_{m=0}^\infty
\frac{\zeta^m}{\sqrt{m!}}|-m,m\rangle.
\label{eq:feqo-vacuum-joint-state}
\end{equation}
负标号只属于电子能量格，光子数 $m$ 始终非负。
两侧基分别正交，所以
\begin{align}
\lambda_m&=\ee^{-\mu}\mu^m/m!,\label{eq:feqo-vacuum-schmidt}\\
\rho_e&=\sum_m\lambda_m|-m\rangle\langle-m|,
\label{eq:feqo-vacuum-reduced-electron}\\
\operatorname{Tr}\rho_e^2&=\ee^{-2\mu}I_0(2\mu),
\label{eq:feqo-vacuum-purity}\\
S_e&=\mu(1-\ln\mu)+\sum_m\lambda_m\ln(m!).
\label{eq:feqo-vacuum-entropy}
\end{align}
第一式的概率和是 $\ee^{-\mu}\sum_m\mu^m/m!=1$。当 $\mu=0$ 时，熵式中的 $\mu\ln\mu$ 按连续极限取零，输出仍为积态。
纯度式由 $I_0(2\mu)=\sum_m\mu^{2m}/(m!)^2$ 得到。
自然对数下，$\mu=0.5,1,4$ 的熵分别为
$0.92764,1.30484,2.08667$；$\mu=1$ 时纯度为 $0.30851$。
大 $\mu$ 时主峰宽度为 $O(\sqrt\mu)$，局部高斯渐近给
$S_e=\tfrac12\ln(2\pi\ee\mu)+O(\mu^{-1})$。
这表示对数增长，并没有在某个有限值饱和。

若以电子损失 $m\hbar\omega$ 为条件，光场精确成为 $|m\rangle$。
无条件光态的数分布虽为 Poisson，密度矩阵却在数态基对角，故
$\langle a\rangle=0$；相干态 $|\alpha\rangle$ 可具有相同数分布却有
$\langle a\rangle=\alpha$。数分布相同并不等于量子态相同。
\end{exbox}

\begin{exbox}[光学记录重叠、纠缠和量子擦除]
对归一化态 $|\Psi\rangle=\sqrt p|0\rangle|u\rangle
+\ee^{\ii\phi}\sqrt{1-p}|1\rangle|v\rangle$，令 $s=\langle v|u\rangle$。
设 $0\le p\le1$，$|u\rangle,|v\rangle$ 为归一化光学记录。求电子相干与纠缠怎样依赖 $s$，再判断读取光学记录能否恢复条件电子相干。

\textbf{解。} 对光求偏迹后，电子的非对角元为
\begin{equation}
(\rho_e)_{01}=\sqrt{p(1-p)}\ee^{-\ii\phi}s,
\quad\operatorname{Tr}\rho_e^2=p^2+(1-p)^2+2p(1-p)|s|^2.
\label{eq:feqo-which-path-overlap}
\end{equation}
约化态两个本征值为
$\lambda_\pm=[1\pm\sqrt{1-4p(1-p)(1-|s|^2)}]/2$，
所以熵可用这两个权重直接计算。
光学记录相同给纯电子态；记录正交且 $p=1/2$ 给熵 $\ln2$。

若反冲与选择规则已受控地隔离联合二态
$|0,1\rangle,|1,0\rangle$，一次半转移可制备
$(|0,1\rangle-\ii|1,0\rangle)/\sqrt2$。
这是一个受限模型，不能在无限无反冲交换链上任意扔掉 $|-1,2\rangle$ 来获得。
读光到 $|\pm\rangle=(|0\rangle\pm|1\rangle)/\sqrt2$，每个结果概率为 $1/2$，
电子条件态分别为 $(|0\rangle\mp\ii|1\rangle)/\sqrt2$。
\begin{equation}
\widetilde\rho_{e|y}=\operatorname{Tr}_{\rm ph}
[(I\otimes P_y)\rho_{e,\rm ph}(I\otimes P_y)],
\quad p_y=\operatorname{Tr}\widetilde\rho_{e|y},
\quad\rho_{e|y}=\widetilde\rho_{e|y}/p_y.
\label{eq:feqo-conditional-electron-state}
\end{equation}
按两个概率加权，条件纯态的相干互相抵消，恢复无条件混态。
恢复条件干涉需要保留结果标签，而不是对两种归一化态不加权求和。
\end{exbox}
